\documentclass{beamer}
\usepackage{graphicx}      % Provides the floating 
\usepackage{listings}
\usepackage{xcolor}

\lstset{
  basicstyle=\ttfamily\footnotesize,
  breaklines=true,
  frame=single
}
%Information to be included in the title page:
\title{Unified Memory Access Profiling}
\author{Sai Coumar}
% \institute{Purdue University}
% \date{2025}
\usetheme{Madrid}

\begin{document}

\frame{\titlepage}

\begin{frame}

\frametitle{Memory Access Profiling Experiments}
\begin{itemize}
    \item A: Host to Device Transfer Baseline  
    \item B: Paging and Migration Behavior
    \item C: Access Pattern Sensitivity
    \item D: Performance under Oversubscription Conditions
\end{itemize}

\end{frame}


\begin{frame}
\frametitle{A: Memory Access Strategies}
\begin{itemize}
    \item Pageable 
    \item Pinned 
    \item Zero Copy
    \item Managed 
\end{itemize}

\end{frame}


\begin{frame}
\frametitle{A: Pageable Memory Access}
\begin{itemize}
    \item Standard Memory Allocation Paradigm
    \item CPU dynamic allocates in system RAM, GPU allocates VRAM, and kernel performs a memory copy
    \item OS overhead increases Host-2-Device (H2D) \& Device-2-Host bandwidth (D2H)
    \item CUDA SDK:
        \begin{itemize}
            \item cudaMalloc()
            \item cudaMemcpy()
        \end{itemize}
\end{itemize}

\end{frame}


\begin{frame}{A: Pinned Memory Access}
\begin{itemize}
    \item Similar paradigm as Pageable Memory Accesses, but pins CPU memory to avoid OS paging in memory
    \item CPU pins data in physical location in system RAM, GPU allocates VRAM, and kernel performs a memory copy
    \item No OS overhead in H2D/D2H memory copy
    \item CUDA SDK:
        \begin{itemize}
            \item cudaMallocHost()
            \item cudaMemcpy()
        \end{itemize}
\end{itemize}
\end{frame}

\begin{frame}{A: Zero-Copy Memory Access}
\begin{itemize}
    \item Memory is allocated in system RAM, never copied entirely into GPU VRAM
    \item GPU Kernel directly accesses memory from system RAM using PCIe
    \item When memory transfer speeds dominate runtime over compute, adds a bottleneck
    \item CUDA SDK:
        \begin{itemize}
            \item cudaHostAlloc()
            \item cudaHostAllocMapped
            \item cudaHostGetDevicePointer()
        \end{itemize}
\end{itemize}
\end{frame}

\begin{frame}{A: Managed Memory Access}
\begin{itemize}
    \item Create a virtual shared memory mapping, move data dynamically on demand
    \item Virtually, ownership bits are changed; Physically depends on the device available 
    \item Discrete GPU will automatically copy memory from system RAM to GPU VRAM while unified architecture SoC will keep everything in the same place physically and only virtually change ownership
     \item CUDA SDK:
        \begin{itemize}
            \item cudaMallocManaged()
        \end{itemize}
\end{itemize}
\end{frame}
\begin{frame}{B: Managed Memory Access Paging and Migration Behavior}
\begin{itemize}
    \item Managed Memory Allocation with no Prefetch
    \item Managed Memory Allocation with Prefetch
    \item Managed Memory Allocation with Thrashing
\end{itemize}
\end{frame}

\begin{frame}{B: Managed Memory Allocation with no Prefetch}
\begin{itemize}
    \item Normal Managed Memory Allocation 
    \item Lazy memory placement performed at runtime
    \item GPU will page fault on first access when data is needed
    \item CUDA SDK:
        \begin{itemize}
            \item cudaMallocManaged()
        \end{itemize}
\end{itemize}
\end{frame}

\begin{frame}{B: Managed Memory Allocation with Prefetch}
\begin{itemize}
    \item Managed Memory Allocation with explicit data migration prior to kernel launch  
    \item Eager memory placement performed prior runtime
    \item Kernel can utilize defined locality for speedup
    \item CUDA SDK:
        \begin{itemize}
            \item cudaMallocManaged()
            \item advisePreferGpu()
        \end{itemize}
\end{itemize}
\end{frame}

\begin{frame}{B: Managed Memory Allocation under Thrashing}
\begin{itemize}
    \item Normal Managed Memory Allocation but with high H2D transfers within the kernel
    \item Lazy memory placement but forced to do many memory allocations 
    \item CUDA SDK:
        \begin{itemize}
            \item cudaMallocManaged()
            \item advisePreferGpu()
        \end{itemize}
\end{itemize}
\end{frame}


\begin{frame}
\frametitle{C: Access Pattern Sensitivity}
\begin{itemize}
    \item Sequential Access 
    \item Sparse Access
    \item Strided Access
\end{itemize}
\end{frame}

\begin{frame}[fragile]{C: Sequential Access}
\begin{itemize}
    \item Iterate over every element in contiguous memory
    \item Best cache reuse conditions
    \item Ideal bandwidth use
\end{itemize}

\begin{lstlisting}[language=C++]
__global__ void k_sequential(float* d, size_t n, float f) {
    size_t i = blockIdx.x * blockDim.x + threadIdx.x;
    if (i < n) d[i] *= f;
}
\end{lstlisting}
\end{frame}

\begin{frame}[fragile]{C: Sparse Access}
\begin{itemize}
    \item Completely random accesses
    \item Poor locality
    \item Triggers many page faults
\end{itemize}

\begin{lstlisting}[language=C++]
__global__ void k_sparse(float* d, size_t n, float f) {
    size_t i = blockIdx.x * blockDim.x + threadIdx.x;
    if (i < n) {
        size_t idx = random(i) % n;
        d[idx] *= f;
    }
}
\end{lstlisting}

\end{frame}

\begin{frame}[fragile]{C: Strided Access}
\begin{itemize}
    \item Make jumps in accesses just large enough to trigger cache misses
    \item Tests sensitivity to non-contiguous accesses
\end{itemize}

\begin{lstlisting}[language=C++]
__global__ void k_strided(float* d, size_t n, int stride, float f) {
    size_t i = blockIdx.x * blockDim.x + threadIdx.x;
    if (i * (size_t)stride < n) d[i * (size_t)stride] *= f;
}
\end{lstlisting}

\end{frame}

\bibliographystyle{IEEEtran}   % or plain, unsrt, apalike
\bibliography{references}
\end{document}
